% Condensed abstract (5-page spine, 2026-10-09). Every number traces to a
% manifested result CSV; see the header comment of each table in paper/tables/.
A decade of work predicts post-earnings-call stock volatility from call
transcripts and audio. We show that the evaluation substrate this literature
shares cannot support its claims: the canonical label and split files have a
zero-day embargo against 30-day targets, place 44--80\% of test calls' companies
in training, and anchor day~0 on the call date although the earnings release
precedes the open for most calls. We release \ecvol-bench, an open benchmark
over FinCall-Surprise and MAEC that repairs this: reconstructed company identity
and dates for a corpus that ships without them, EDGAR-measured release
timestamps, targets under both horizon conventions, embargoed and
\tickerdisjoint{} splits, econometric baselines, and an identity-control suite
that runs on any prediction file. Under these controls, modern frozen text and
audio representations add nothing to past volatility: adding either to a
past-volatility ridge lowers \Rtwo{} in all 16 split$\times$horizon cells, text
embeddings identify the company at 89.5\% probe accuracy, audio-only heads do no
better than a featureless intercept, and fusion beats HAR-RV in one horizon. Re-run on their
own data, the field's text models lose 12--82\% MSE when companies are held out
while price-based models do not move. Nothing degrades when the frozen pipeline
is scored on 475 calls held after every encoder's knowledge cutoff. All code,
metadata, splits, targets, features and result tables are released.
